% !TeX root = ../papers/nc-p3.tex
\section{四つの例外倍率の重根上の全次数正則性}
\label{sec:exceptional-ramification}
\begin{theorem}\label{thm:exceptional-ramification-as}
$z\ne0,\pm1$, $D_0\ne0$として$\lambda=\pm1,\pm1/z$を取る。
各判別四次は滑らかな二次曲面を二重に数えたものである。この二次曲面上の
可逆な窓が標準Koszul左右AS正則となる必要十分条件は、$P,Q$がともに真の二次式
であることである。その場合の完成は一意である。
\end{theorem}
\begin{proof}
$G=(a+d)(zb+c)-(z+1)(ad+bc)$と置く。正の二倍率$1,1/z$での$F$は
それぞれ$(z-1)^2G^2,(z-1)^2G^2/z^2$で、負の倍率では$b,c$の符号を変える。
$G$のHesse行列式は$(z^2-1)^2$である。真の重根の完成は
定理\ref{thm:klein-other-pairings}の表による。$pq=0$の有限解では、
その側の線型係数も判別式の零から消え、定数項も零になるので、
定理\ref{thm:all-zero-pairing-defect}, \ref{thm:single-zero-pairings}で除かれる。

まず$\lambda=1$とし、$A=A_0=A_1$とする。実際には$G=0$より広く、
次の八差が全て非零な全ての$A$で標準正則性を示す。
\[
 a-b,\quad a-c,\quad a-zb,\quad za-c,\quad
 b-d,\quad c-d,\quad zb-d,\quad c-zd .
\]
$G=0$では$x=y=D_0$なので真の重根なら$n=pD_0^2$と$o=qD_0^2$は非零である。
各差を零にすると、$G$の制限は$(a-zd)(za-d)(b-c)(z^2b-c)$の制限を割る。
例えば$a=b$では$G=(b-c)(zb-d)$である。従ってこの八差は全て非零となる。

両対称表示で反対切断を文字次数に分ける。四行二列、二行四列の二ブロックと、
行列式$z-1,1-z$の二行二列ブロック二つを得る。
長方形ブロックを両方とも四行二列に向けて行対を$01,02,03,12,13,23$の順に取る。
最初と最後の小行列は単元倍で順に$(a-d,zb-c)$と$(zb-c,a-d)$である。
中の四つはすべて上の八差のうち二つの積を$D_0$で割ったものとなる。
従って反対階数は8で、左右の二次核は四つの平方だけからなる二次鉛筆である。
行対の小行列は核二平面の相補小行列に対応する。
$a-d,zb-c$のいずれも零でなければ、各鉛筆の四つの列方向は相異なり、両曲線は滑らかである。
片方だけが零なら、両曲線とも横断する二平面二次曲線の環である。
実際、一対の列が比例するのでその二列を消す一式は残る二変数の異なる二直線因子に分かれ、
他方の式は各平面に非退化な対角二次曲線を切る。二対とも比例すれば、
独立な二組それぞれの二変数二次式が分解し、両曲線とも四直線の環となる。
これらの曲線の埋込み次数はそれぞれ$(2,2)$または$(1,1,1,1)$である。

六次元の関係空間に純テンソルがないことを確かめる。
文字を低い二文字と高い二文字に分けると、関係行列の右上と左下は
\[
 \begin{pmatrix}r&t\\t&r\end{pmatrix},\qquad
 \begin{pmatrix}r&t\\t&zr\end{pmatrix}.
\]
両行列式が零なら$z\ne1$より$r=t=0$である。残るブロック対角行列が階数一なら
対角ブロックの片方が零であるが、各厚い文字次数の関係二平面は低い二座標にも
高い二座標にも同型に射影するので、全体が零である。
従って滑らかな場合には定理\ref{thm:smooth-apolar-complete}が適用できる。

多角形の場合に必要な両側評価の階数を示す。零点対の低い二座標を$u,s$、
高い二座標を$v,t$とする。厚い四式は
\[
 \binom{u_0s_0}{u_1s_1}=L\binom{v_0t_0}{v_1t_1},\qquad
 \binom{u_0s_1}{u_1s_0}=K\binom{v_0t_1}{v_1t_0},\qquad
 L=-A^{-t}Y^t,\quad K=-A^{-1}X
\]
で、$X=\left(\begin{smallmatrix}1&1\\1&z\end{smallmatrix}\right)$,
$Y=\left(\begin{smallmatrix}1&1\\z&1\end{smallmatrix}\right)$である。
全八成分が非零で、$\det K=-\det L=(z-1)/D_0\ne0$である。
$(u,v)$を固定した評価の核が二次元以上と仮定する。
$u$または$v$が零なら厚い式が対応する二座標を消し、薄い式がさらに一式を課すので不可能である。
どちらも非零なら、核から$s$への射影も$t$への射影も単射である。
従って$t$は自由に動き、$s$も二次元を張る。上の四式の成分は全て非零なので、
$v$の一座標が零なら、例えば$v_1=0$で
$(u_0s_0)(u_1s_1)-(u_0s_1)(u_1s_0)$の右辺の$t_0^2$係数が非零となり矛盾する。
同様に全ての$u_i,v_i$が非零となる。
$t_0,t_1$の係数を比較し、$u_1/u_0$を二通りに求めると
$K_{11}K_{00}=L_{00}L_{11}$かつ$K_{10}K_{01}=L_{01}L_{10}$となる。
これは行列式の符号に反する。右評価では$K$の行と列を同時に交換するだけなので同様である。

滑らかな楕円鉛筆からこの窓への特殊化を取れる。実際、八差を非零に保った係数空間は
既約で、一般点の二鉛筆は滑らかであり、上でその正則表示を得た。
多角形への特殊化で、定理\ref{thm:two-polygon-specialization}の証明の共通の環を取る。
多重次数が異なる174型中172型では、同証明が実際の関係の片側評価に
二次元以上の核を与える。残る二型はどちらも次数3の成分を必要とし、
ここで得た埋込み次数に存在しない。従って両束の多重次数は等しく、同証明により
八次元の積の上の二次元鉛筆が得られる。基点は純テンソルを与えるため存在しない。
多角形鉛筆の全次数定理で正則となる。

$\lambda=1/z$の表の窓は$A'=\left(\begin{smallmatrix}zb&a\\zd&c\end{smallmatrix}\right)$を持つ
$\lambda=1$の窓と同型である。具体的に第三・第四位置の文字を各々1だけ平行移動し、
成分に$(1,1,1/z,1)_k(z,1,1,1)_l$を掛ける。$D'_0=-zD_0$,
$G(A')=zG(A)$で、$no$も$z^2$倍となる。負の二倍率は既証の可逆な
コサイクル捻りで正の二倍率へ戻り、$b,c$がともに反転する。
\path{computations/verify_exceptional_ramification.py}が四つの判別四次、上の小行列、
八差への制限、正倍率の基底変更の全256成分を検算する。
\end{proof}
